\documentclass{beamer}

\usepackage[utf8]{inputenc}
\input{../helper}

\title{Introdução à Computação \\ tipos abstratos, pilhas}
\author{Alexandre Rademaker}
\date{}

\begin{document}

\begin{frame}
 \maketitle
\end{frame}

\begin{frame}[allowframebreaks]{tipos abstratos de dados e estruturas de dados}

  Um \emph{tipo abstrato de dados} (ADT) é um modelo para tipos de
  dados definido por seu comportamento do ponto de vista de um usuário
  dos dados, especificamente em termos de valores possíveis, operações
  possíveis em dados de esse tipo e o comportamento dessas operações.

  Uma \emph{estrutura de dados} é um formato especializado para
  organizar, processar, recuperar e armazenar dados. Existem vários
  tipos (básicos e avançados) de estruturas de dados, geralmente
  projetadas para finalidades específicas como garantir uma dada
  complexidade para alguma operação.

  \framebreak

  O tipo \texttt{int} de C é um tipo de dados que implementa o tipo de
  dados abstrato \emph{inteiro}.

  Inteiros podem ser representados como números binários, na maioria
  das vezes como \textbf{complemento de dois}, mas podem ser
  \textbf{decimais codificados em binário} ou em \textbf{complemento
    de um}.

  Usamos inteiros e as operações sobre inteiros sem nos preocupar com
  a implementação binária dos inteiros e como as operações são
  implementadas em termos da codificação binária.

  \framebreak

  O tipo abstrato de dados \emph{lista} (ou sequência) representa uma
  coleção de valores ordenada.

  Sequências podem ser implementadas como \textbf{arrays} ou
  \textbf{listas encadeadas}.

  As operações mais comuns são: incluir um elemento, acessar o nésimo
  elemento, remover um elemento, localizar um elemento, concatenar
  duas sequências, ordenar, etc.
  
\end{frame}


\begin{frame}[allowframebreaks,fragile]{pilhas}

  Em uma pilha (stack), a manipulação dos elementos é realizada em
  apenas uma das extremidades, chamada de topo, em oposição a outra
  extremidade, chamada de base.

  \begin{center}
    \includegraphics[width=.6\textwidth]{stack.png}
  \end{center}

  \framebreak

  Pilha um tipo abstrato de dados do tipo ``Last-In, First Out''
  (LIFO), ou seja ``o último que entra é o primeiro que sai''.

  Pilhas podem ser implementadas em estruturas de dados concretas como
  arrays e listas encadeadas.
  
  \framebreak

  As operações básicas de uma pilha são:

  \begin{description}
  \item[push] insere um novo elemento no topo
  \item[pop] remove o elemento do topo
  \item[peek] retorna o elemento no topo sem remover da pilha
  \item[isEmpty] testa se a pilha está vazia.
  \item[size] retorna o tamanho da pilha, quantos elementos ela
    contém.
  \end{description}

  Todas podem ter complexidade $O(1)$ quando pilhas são implementadas
  usando arrays ou listas encadeadas.

  \framebreak

  Pilhas podem ser usadas para reverter uma sequência (lista encadeada
  ou array) e ordenar uma sequência.

  Verificar se uma expressão aritmética tem os parênteses
  balanceados. Por exemplo, a expressão $(2 + (8 / y))$ e a expressão
  $(2 + (8 / y)$.

  Implementar um mecanismo de 'desfazer' operações.

  Pilhas são ainda usadas para implementar buscas em grafos, conversão
  de expressões aritméticas em notação infixada ($2 + 3$) para
  pré-fixada ($+ 2 3$) ou pós-fixada ($2 3 +$).
\end{frame}


\begin{frame}[allowframebreaks,fragile]{pilhas usando array}

  \begin{center}
    \includegraphics[width=\textwidth]{stack-as-array.png}
  \end{center}

  \framebreak

  \begin{lstlisting}[style=normalc]
    // maximum size of the stack
    #define MAX_SIZE 100 

    typedef struct stack {
      int stack[MAX_SIZE];
      int top; 
    } t_stack;
    
  \end{lstlisting}

  Implementando com array! Desvantagem, alocação estática com limite
  pré-definido de elementos. Vantagem, simplicidade.

  \framebreak

  \begin{lstlisting}[style=normalc,caption=src/stack-array.c]
    t_stack* make_stack(){
      t_stack *t = malloc(sizeof(t_stack));
      t->top = -1;
      return t;
    }
    
    void push(t_stack *s, int data) {
      assert(s->top < MAX_SIZE - 1);
      s->stack[++s->top] = data;
    }
    
    int pop(t_stack *s) {
      assert(s->top > -1);
      return s->stack[s->top--];
    }
  \end{lstlisting}

\end{frame}


\begin{frame}[allowframebreaks,fragile]{pilhas usando listas encadeadas}

  \begin{center}
    \includegraphics[width=\textwidth]{stack-as-list.png}
  \end{center}

  \framebreak

  \begin{lstlisting}[style=normalc]
    typedef struct node {
      int number;
      struct node *next;
    } t_node;
  \end{lstlisting}

  Mesma \texttt{struct} declarada anteriormente para listas
  encadeadas. Nenhuma inicialização é necessária.

  \framebreak

  \begin{lstlisting}[style=normalc,caption=src/stack-list.c]
    int pop(t_node **p) {
      assert(*p != NULL);
      int r = (*p)->number;
      t_node *t = *p; *p = (*p)->next; free(t);
      return r;
    }
    
    void push(t_node **p, int e) {
      t_node* n = malloc(sizeof(t_node));
      assert(n != NULL); n->number = e; n->next = *p;
      *p = n;
    }

    int peek(t_node *p) {
      assert(p != NULL); return p->number;
    }
  \end{lstlisting}

  \framebreak

  \begin{lstlisting}[style=normalc]
    int main(int argc, char *argv[]) {
      t_node *s = NULL;
      push(&s, 1); 
      push(&s, 2); 
      push(&s, 3);
      
      printf("popped %d \n", pop(&s));
    }
  \end{lstlisting}

  \texttt{push} e \texttt{pop} alteram a pilha, por isso passamos o
  endereço de \texttt{s} (endereço da variável que é um endereço de um
  \texttt{t\_node}) não o valor de \texttt{s} (endereço de um
  \texttt{t\_node}).

  \texttt{push} e \texttt{pop} alteram o valor de \texttt{s}, declarada
  na \texttt{main}.

  \framebreak

  \begin{center}
    \includegraphics[width=\textwidth]{double-pointer.png}
  \end{center}

\end{frame}


\end{document}

%%% Local Variables:
%%% mode: latex
%%% TeX-master: t
%%% End:
